%----------------------------------------------------------------------
\section{Discussion}\label{sec:discussion}
%----------------------------------------------------------------------

\subsection{Summary of results}

The following table summarizes the correspondence
between [DAF]'s foundational machinery and the
scheduling machinery developed here.

\medskip
\begin{center}
\begin{tabular}{l|l|l}
\textbf{[DAF] construct} & \textbf{This paper's analogue} & \textbf{Status} \\
\hline
Discriminator $d$ & Atom $a \in \mathcal{A}$ & Definition \\
$\kappa$-evolution & Chart item insertion & Theorem \ref{thm:soundness} \\
Membership profile $(\tau, \sigma)$ & Semantic attributes $(E, r)$ & Definition \\
Four-cell structure & Encoding lattice $\mathcal{E}$ & Definition \ref{def:lambda}ff. \\
Residue $R(S)$ & Cost-vector slack $B - \bar{c}$ & Definition \ref{def:cost-vector} \\
Extensionality & Plan-equivalence (Definition \ref{def:exec-operationalist}) & Proposition \ref{prop:plan-trans} \\
Operationalist & Execution-operationalist & Definition \ref{def:exec-operationalist} \\
Gluing axiom & Cost-bounded completer & Definition \ref{def:cost-bounded} \\
Sheafification & Chart closure & Theorem \ref{thm:earley-complete} \\
SPPF & Plan portfolio & Definition \ref{def:portfolio} \\
Fano closure & LCM density & Proposition \ref{prop:lcm-density} \\
Yoneda embedding & Cost-signature equivalence & Proposition \ref{prop:plan-trans} \\
Chain complex $\partial \circ \partial = 0$ & Grade-drop chain under $\partial$ in RM & Section \ref{sec:encoding} \\
\end{tabular}
\end{center}

Each row is a functorial correspondence: the
structural role [DAF] assigns to the left-hand
concept is played, in this paper, by the right-hand
concept. The parallel is close enough that both
papers can be read as instances of a single
underlying pattern: sheafification applied to the
derivation DAG of an open-world computation, with
the sheaf's sections being local plans or local
set-membership claims, and the gluing axiom being
satisfied by either cost-bounded parsing or
$\GF(2)$-linear constraint solving.

\subsection{Relationship to existing work}

The components are drawn from well-established
literatures. The novelty lies in the specific
synthesis.

\begin{itemize}[nosep]
  \item \textbf{Earley parsing with attributes}.
    Knuth (1968), Earley (1970), Aycock \& Horspool
    (2002). The attribute-grammar extension of
    Earley is classical; the specialization of its
    attribute domain to a multi-dimensional cost
    lattice is well-known in weighted CFG parsing
    (Chiang et al., Huang, Lopez) for natural-language
    translation applications.
  \item \textbf{Lagrangian relaxation for CFG parsing}.
    Chiang, Knight, Haghighi (2008) introduced
    Lagrangian techniques for constrained
    decoding; Huang (2008) extended them to
    $k$-best parsing. Our contribution is to
    multi-resource constraints, where the dual
    variable is itself vector-valued.
  \item \textbf{Pareto-A$^*$ search}. Mandow \&
    P\'{e}rez-de-la-Cruz (2010) give the
    general-purpose algorithm for multi-objective
    shortest paths; we import their framework
    essentially unchanged.
  \item \textbf{Block-diagonal matrix fusion}.
    Standard in BLAS implementations for
    block-sparse problems; the scheduling
    application to heterogeneous atoms sharing a
    SIMT register is a specialization not found in
    the BLAS literature.
  \item \textbf{GF(2) matrix methods}. Valiant
    (1975) reduces CFG parsing to matrix
    multiplication; the converse, recognizing a
    parse as a GF(2)-linear algebraic computation,
    is implicit in our grammar but not exploited for
    the parsing algorithm itself (we use Earley,
    not matmul-based methods).
  \item \textbf{Kernel fusion and software
    pipelining}. The compiler-optimization
    literature (Muchnick 1997; Allen \& Kennedy
    2001) gives the algorithmic foundation for
    kernel fusion; our \texttt{fused\_matrix\_apply}
    is a specific instantiation of that idea for
    $\GF(2)$-linear atoms over masked registers.
\end{itemize}

The novelty is the \emph{composition}: a grammar
whose productions are GPU scheduling
transformations, whose attributes are hardware
resource vectors, whose parser is Earley with
cost-bounded completion, whose output is a Pareto
plan portfolio. Each piece is individually classical;
their synthesis as an integrated compiler for the
CSTZ algebra is new.

\subsection{Practical limitations}

\paragraph{Calibration burden.} The hardware profile
$\pi$ is the locus of device-specific information.
Calibrating $\pi$ from published SM specifications
is plausible but non-trivial; calibrating from
measured benchmarks on a specific device requires a
benchmark sweep across the atom inventory. This is
the principal engineering investment required to
bring the planner from design to operation.
(\S1 item S1, Appendix \ref{app:kernels} hints at
the measurement protocol.)

\paragraph{Grammar completeness.} The atom inventory
of Section \ref{sec:atoms} covers the CSTZ algebra
as developed in [DAF]'s Sections 3--9. Extensions to
downstream frameworks --- adjoint functor lifting,
higher-category morphism composition, sheaf
restriction maps --- would add atoms. Each new
atom requires a cost signature and a property test;
the grammar structure accommodates arbitrary
additions, but completeness against any particular
downstream framework is the user's responsibility.

\paragraph{Non-linear atoms.} The XOR-bridge and
fused-matrix mechanics of Sections
\ref{sec:fusion}--\ref{sec:time-sharing} apply only
to $\GF(2)$-linear atoms. Popcount, comparison,
and other reductions force register spills. Whether
a broader class of atoms admits a ``semi-reversible''
extension of the XOR-bridge is the subject of
Conjecture \ref{conj:nonlinear}.

\paragraph{Dynamic shapes.} The planner presupposes
static $n$ (exterior dimension). Workloads that
require runtime $n$-variability would need to be
split into a family of plans, one per $n$, with
dispatch to the correct plan selected at kernel
invocation. This is expressible in the grammar but
adds complexity to the parsing stage and is deferred.

\paragraph{LP integrality gap.} The Lagrangian
relaxation of Section \ref{sec:lagrangian} has a
provably-zero gap in the affine-cost case, but the
general worst-case gap is only empirically bounded.
A tight theoretical bound is listed as open (X3).

\subsection{Future directions}

\paragraph{Grammar-level verification in Agda.}
The predecessor [DAF] is mechanized in Agda
(\texttt{agda/CSTZ/}); the planner grammar
could similarly be formalized, with each
production type-checked against an Agda
specification of the underlying $\GF(2)$ atom
semantics. This would give a mechanical proof
of Theorem \ref{thm:soundness} at the grammar
level.

\paragraph{Streaming planner for ObservationState.}
[DAF]'s classification subsystem emits an
ObservationState that grows monotonically as
new classifier observations are ingested.
Proposition \ref{prop:incremental} lets the
planner update in place; a streaming integration
between the Python ObservationState
(\texttt{src/cstz/observe.py}) and the planner's
chart would give real-time plan updates as new
observations arrive.

\paragraph{Cost-profile calibration framework.}
A benchmark suite that measures per-atom resource
consumption on a target device, then fits a
hardware profile $\pi$, would close item S1 of
the reading guide. This is essentially the same
shape as a BLAS autotuning loop.

\paragraph{Multi-device scheduling.} The current
budget vector is per-device. Extending the grammar
to a multi-device cluster (data-parallel or
model-parallel) would add a communication cost
dimension and would require a production rule for
cross-device atoms. This is a natural extension
but not attempted here.

\subsection{The operationalist through-line}

Both [DAF] and this paper rest on an operationalist
commitment: that a distinction which cannot be
witnessed is not a valid distinction. [DAF]
applies this to the metaphysics of discrimination
(``if no discriminator separates $a$ and $b$ under
any $\kappa$-evolution, then $a = b$''). We apply
it to the metaphysics of execution (``if no
measurement separates $P_1$ and $P_2$ on the
target hardware, then $P_1$ and $P_2$ are
execution-equivalent'').

The commitment is philosophically the same. In
both cases, identifying two entities that would
otherwise be distinct is an engineering act, not a
forced mathematical conclusion; in both cases, the
engineering motivation is that the framework that
admits un-witnessable distinctions cannot be
implemented, tested, or falsified. The planner's
portfolio output, which retains alternative plans
whose costs are measurement-indistinguishable, is
the execution-level manifestation of the same
principle that licenses [DAF]'s sheafification to
treat $\kappa$-equivalence as identity.

\subsection{Governance as parity check}

The verification map (Section \ref{sec:reading-guide})
is not an editorial nicety: it is [DAF]'s XOR
signature $\chi = \tau \oplus \sigma$ applied to the
paper itself. Each claim is probed by two
independent discriminators --- \textsf{proven} and
\textsf{witnessed} --- whose XOR gives the claim's
status:
\begin{center}
\begin{tabular}{c|c|c}
\textsf{proven} & \textsf{witnessed} & mark \\
\hline
1 & 1 & \cmark\ (``verified'': argument complete AND
    witness exists) \\
1 & 0 & \gmark\ (``gap'': argument complete, specific
    witness missing) \\
0 & 1 & --- (inadmissible: a witnessed but
    unproven claim is not a claim) \\
0 & 0 & \omark\ (``open'': framework declines to
    prove, not a defect)
\end{tabular}
\end{center}

The \gmark\ mark is exactly the four-cell structure's
\emph{residue}: a signature with $\chi = 0$ that
signals a discrimination is available but not yet
performed. [DAF]'s residue analysis prescribes what to
do with such signatures: treat the residue as a
discriminator-articulation request and $\kappa$-evolve
the regime to resolve it.

Applying this to the paper: each \gmark\ in the
verification map is paired with a named gap item
(G1--G6) that specifies \emph{exactly} which
$\kappa$-evolution would convert the residue to a
\cmark. The framework is not just self-diagnostic;
it is \emph{self-prescriptive}. The following table
records the parity check's output:

\medskip
\begin{center}
\footnotesize
\begin{tabular}{l|c|l|l}
\textbf{Claim} & \textbf{Mark} & \textbf{Residue} & \textbf{Required $\kappa$-evolution} \\
\hline
Complexity (\ref{thm:complexity}) & \gmark &
  $|\mathrm{Pareto}|$ unbounded WC & G3: measure frontier on real workloads \\
Heuristic monotone (\ref{prop:heuristic-monotone}) & \gmark &
  $h^*$ near-trivial & G6: measure pruning rate; strengthen if needed \\
Parse-packing (\ref{thm:parse-packing}) & \gmark &
  Def \ref{def:atom-equiv} underdet.\ & G4: algorithmic decidability of atom-seq.\ equiv.\ \\
Reg-time-space (\ref{thm:reg-bound}) & \gmark &
  Rectangular abstraction & G5: stepped-profile generalization \\
Lagrangian (\ref{thm:lagrangian}) & \omark &
  Integer gap bound & X3: prove tight bound (or accept open) \\
All \textsf{H}-class \cmark & (conditional) &
  Uncalibrated $\pi$ & G1: fit profile for ref.\ GPU \\
\multicolumn{2}{l}{--- operational validity} &
  No parser binary & G2: implement Earley chart parser \\
\end{tabular}
\end{center}

\medskip

Pareto dominance (Section \ref{sec:encoding} onward)
is itself an instance of the parity-check mechanism.
For two plans $P, P'$, the \emph{admissibility
XOR-signature} at budget $B$ is
\begin{equation*}
  \chi_B(P, P') = \delta_B(P) \oplus \delta_B(P').
\end{equation*}
$\chi_B = 0$ when both are equally admissible under
$B$ (both fit or neither fits); $\chi_B = 1$ when
the admissibility regime distinguishes them. Plan
$P'$ operationalist-dominates $P$ iff $\chi_B = 0$
for every $B \in A(P)$ (every budget that admits $P$
also admits $P'$) and $\chi_B = 1$ for some
$B \in A(P') \setminus A(P)$ (at least one budget
admits $P'$ but not $P$). The witness-budgets in the
asymmetric difference $A(P') \setminus A(P)$ are the
residues; they index exactly the dimensions on which
$c(P) - c(P') \geq 0$ has strict-positive support.
This is the admissibility regime's XOR signature,
applied reflexively to the paper's exclusion rule.

The framework's self-regulation applied recursively:
a reader uncomfortable with an \emph{un}derived
exclusion rule (as was rightly the case in an
earlier draft that imported Pareto dominance
wholesale) has detected an operationalist-unwitnessed
claim. The remedy (Propositions
\ref{prop:op-equals-pareto} and
\ref{prop:frontier-char}) supplies the witness that
converts the $\gmark$ to $\cmark$ without moving any
cost vector or phase boundary; only the
\emph{derivation} changes, from imported MOO
convention to algebraic consequence of
exec-operationalism on the admissibility regime.

\medskip

Every residue has a named discriminator. A reader who
feels the rigor-differential between [DAF]'s
\textsf{A}-class claims (backed by mechanized Agda)
and this paper's \textsf{H}-class claims (backed by
``a suitable profile'') has correctly detected
residues G1 and G2. The framework's response is not
to explain away the discomfort but to act on it:
close G1 and G2, and the proportion of \cmark\ marks
on \textsf{H}-class claims equals the proportion on
\textsf{A}-class claims automatically, without
changing any definition or proof.

In this sense, the paper is a properly-labeled
intermediate state of a program that [DAF]
initiated. The relationship between [DAF] and this
paper is not ``predecessor--successor'' in the sense
of two independent works; it is $\kappa_t$ and
$\kappa_{t+1}$ in the same regime, where the
evolution from one to the next is the articulation
of the execution-operationalist axiom
(Definition \ref{def:exec-operationalist}) and the
corresponding \textsf{H} / \textsf{R} / \textsf{F}
axiom classes. The next evolution, $\kappa_{t+2}$,
is the paper that substitutes measured numbers for
MODEL-4's symbolic parameters and promotes the
\gmark\ marks in the table above to \cmark\ by
witnessing.

The governance does not ask the reader to suspend
discomfort; it asks the reader to let discomfort
drive the next $\kappa$-evolution.

\subsection{Conclusion}

We have extended the discrimination framework of
[DAF] from mathematical foundations to execution
engineering. A weighted attribute grammar captures
the admissible schedules; an Earley chart parser
with cost-bounded completion selects among them;
Lagrangian duality over a four-dimensional resource
vector gives the constraint-enforcement mechanism;
Pareto portfolios retain non-dominated alternatives;
ambiguous-parse packing exploits the batch coarseness
that LCM tile alignment enforces to run alternatives
concurrently at no extra dispatch cost. The two
compound productions --- spatial MIMD fusion via
block-diagonal masked matrix, and temporal
multiplexing via XOR-bridged register sharing ---
are opposite poles of the fusion-vs-sequencing
tradeoff, and their nesting is what the planner's
search is searching over.

The one genuinely new foundational commitment is
the execution-operationalist axiom
(Definition \ref{def:exec-operationalist}), which
declares plan-equivalence to be the coarsest
relation consistent with observable cost behavior
on stated hardware. Everything else is a lifting
of [DAF]'s algebraic machinery into an execution
model, an operation that---like sheafification over
a site---produces a substantially richer output
than any single derivation would admit.
